\subsubsection{Método}\label{sec:construccion-metodo}

La fase de Construcción se desarrollará mediante un enfoque ágil e iterativo adaptado de la meta de \textit{Construcción} del marco \textit{Disciplined Agile Delivery} (DAD). Esta adaptación metodológica se justifica en tres pilares fundamentales:

\begin{itemize}
    \item \textbf{Enfoque en calidad integral desde el inicio:} DAD promueve la verificación continua y el aseguramiento de la calidad en cada incremento funcional, garantizando que el motor de simulación, la orquestación en HPC y los componentes espaciales cumplan con los estándares de ingeniería de software definidos antes del cierre del sistema.
    \item \textbf{Desarrollo incremental guiado por la investigación:} Permite descomponer la complejidad de la simulación híbrida de incendios en incrementos funcionales y evaluables. Cada iteración arrojará un resultado identificable —sea un módulo ejecutable de software o la evidencia técnica para fijar una configuración, un límite de datos o una condición de consulta— evitando suponer que toda investigación debe generar inmediatamente una función de usuario y evitando reservar la verificación para la entrega final.
    \item \textbf{Gestión rigurosa de entregables y gobernanza:} Facilita el seguimiento disciplinado tanto del código y artefactos reproducibles (paquetes E2 a E6) como de la documentación de calidad y certificación requerida por la universidad y los estándares internacionales adoptados.
\end{itemize}

Cada incremento se implementará, comprobará y revisará sistemáticamente junto con el equipo y el director de tesis, conservando el control de versiones de los artefactos evaluados para asegurar la trazabilidad técnica y científica del proceso.